% Authoritative exercise source. Preserve all derivations and diagrams.

\ExerciseStart{G4.01}{Ricostruire l'interferenza in Deutsch–Jozsa e Bernstein–Vazirani}{Guida esplicita}

\BookHeading{Problema, dati e obiettivo}

Per Deutsch–Jozsa su tre bit confrontare la funzione costante \(f=1\) con \(f(x_0,x_1,x_2)=x_0x_1\oplus x_2\). Per Bernstein–Vazirani usare \(g(x)=x_0\oplus x_2\), associata alla stringa 101. In entrambi i casi preparare l’ausiliario in \(\ket-\) e svolgere tutta l’interferenza.

\ExTheory{\ExRef{c4-deutsch-jozsa}{4.01}; \ExRef{c4-bernstein-vazirani}{4.02}}{L’oracolo reversibile scrive \(y\oplus f(x)\); su \(\ket-\) produce il segno \((-1)^{f(x)}\). La trasformata Hadamard soddisfa \(H^{\otimes3}\ket x=2^{-3/2}\sum_y(-1)^{x\cdot y}\ket y\), dove il prodotto scalare è calcolato modulo due.}

\ExPrediction{La ricombinazione somma segmenti reali di segno opposto. Nel caso bilanciato il coefficiente di 000 deve cancellarsi; la forma affine di Bernstein–Vazirani dovrebbe concentrare l’intera ampiezza sulla stringa cercata.}

\ExFigure{\begin{tikzpicture}[>=Stealth,exlabel]
\node[exbox] (n0) at (0.0,0.0) {prepara $\ket+^{\otimes3}\ket-$};\node[exbox] (n1) at (3.6,0.0) {applica i segni};\draw[exother] (n0)--(n1);\node[exbox] (n2) at (7.2,0.0) {ricombina con $H^{\otimes3}$};\draw[exother] (n1)--(n2);
\end{tikzpicture}}{Due algoritmi condividono la preparazione, mentre la struttura della funzione determina le cancellazioni.}


\begin{ExTasks}

\item Verificare l’equilibrio della prima funzione.

\item Scrivere lo stato dopo il ritorno di fase.

\item Calcolare tutte le otto ampiezze di Deutsch–Jozsa.

\item Confrontare la funzione costante.

\item Fattorizzare la somma di Bernstein–Vazirani sui tre bit.

\item Calcolare la distribuzione e discutere un termine affine costante.

\end{ExTasks}

\ExSolution

\ExStep{Costruire la tabella}

Per \(x=000,001,010,011,100,101,110,111\), i valori di f sono \(0,1,0,1,0,1,1,0\). Ci sono quattro zeri e quattro uni, quindi la promessa di bilanciamento è soddisfatta.

\ExStep{Preparare i segni}

\ExEq{\ket{\Psi}&=\frac1{\sqrt8}\sum_{x=0}^{7}(-1)^{f(x)}\ket x\otimes\ket-,\\
A(y)&=\frac18\sum_{x=0}^{7}(-1)^{f(x)+x\cdot y}.}
Le due normalizzazioni \(1/\sqrt8\) producono il fattore \(1/8\).

\ExStep{Svolgere le otto somme}

\ExEq{A(000)&=(1-1+1-1+1-1-1+1)/8=0\\
A(001)&=(1+1+1+1+1+1-1-1)/8=0.5\\
A(010)&=(1-1-1+1+1-1+1-1)/8=0\\
A(011)&=(1+1-1-1+1+1+1+1)/8=0.5\\
A(100)&=(1-1+1-1-1+1+1-1)/8=0\\
A(101)&=(1+1+1+1-1-1+1+1)/8=0.5\\
A(110)&=(1-1-1+1-1+1-1+1)/8=0\\
A(111)&=(1+1-1-1-1-1-1-1)/8=-0.5}
Gli esiti 001, 011, 101 e 111 hanno probabilità \(1/4\); gli altri hanno probabilità zero. In particolare l’uscita 000 è impossibile.

\ExStep{Confrontare la costante}

Per \(f=1\) ogni segno è \(-1\). Per y=000, \(A(000)=(-1-\cdots-1)/8=-1\). Per ogni y non nullo, le otto parole x contengono quattro prodotti \(x\cdot y=0\) e quattro uguali a uno: la somma è zero. La lettura restituisce 000 con certezza.

\ExStep{Ricostruire la stringa affine}

Per \(g(x)=s\cdot x\), con \(s=101\),
\ExEq{B(y)&=\frac18\sum_x(-1)^{(s\oplus y)\cdot x}\\
&=\prod_{j=0}^{2}\frac{1+(-1)^{s_j\oplus y_j}}2.}
Se \(y_j=s_j\), il fattore j vale \((1+1)/2=1\); se differiscono, vale \((1-1)/2=0\). Quindi \(B(101)=1\) e le altre sette ampiezze sono zero.

\ExStep{Controllare termine affine e risorsa}

Per \(g(x)=s\cdot x\oplus1\), tutte le ampiezze acquistano il fattore \(-1\): la lettura rimane 101. Ciascun circuito usa una interrogazione coerente dell’oracolo, oltre alle porte di preparazione e lettura. Il conteggio è distinto dal costo di costruzione dell’oracolo.

\par\noindent\begin{minipage}{\linewidth}
\ExFigure{\begin{tikzpicture}[>=Stealth,exlabel]
\draw[exaxis] (-.6,0)--(4.4,0) node[right] {esito};\draw[exaxis] (-.6,0)--(-.6,3.6) node[above] {$p$};\node[left] at (-.6,3) {$ 0.3 $};\draw[Ink!15,dashed] (-.6,3)--(3.9000000000000004,3);\filldraw[fill=Accent!25,draw=Accent] (0.0,0) rectangle (0.65,2.5);\node[below] at (0.325,0) {$ 001 $};\node[above] at (0.325,2.5) {$ 0.25 $};\filldraw[fill=Accent!25,draw=Accent] (1.1,0) rectangle (1.75,2.5);\node[below] at (1.425,0) {$ 011 $};\node[above] at (1.425,2.5) {$ 0.25 $};\filldraw[fill=Accent!25,draw=Accent] (2.2,0) rectangle (2.85,2.5);\node[below] at (2.5250000000000004,0) {$ 101 $};\node[above] at (2.5250000000000004,2.5) {$ 0.25 $};\filldraw[fill=Accent!25,draw=Accent] (3.3000000000000003,0) rectangle (3.95,2.5);\node[below] at (3.6250000000000004,0) {$ 111 $};\node[above] at (3.6250000000000004,2.5) {$ 0.25 $};
\end{tikzpicture}}{Esiti di Deutsch–Jozsa per la funzione bilanciata non affine; Bernstein–Vazirani concentra invece tutto il peso su 101.}


\ExInterpret{La distinzione fra funzione costante e bilanciata riguarda il coefficiente 000. L’ipotesi affine più forte permette di recuperare un’intera stringa mediante la stessa somma di caratteri binari.}

\ExLab{Nelle sezioni 4.01 e 4.02 scegliere rispettivamente «Bilanciata non affine su tre bit» e «Stringa 101». Confrontare i segni prima della ricombinazione, poi le distribuzioni finali. Usare anche la variante affine con termine costante per osservare la fase globale.}


\end{minipage}
